% Options for packages loaded elsewhere
% Options for packages loaded elsewhere
\PassOptionsToPackage{unicode}{hyperref}
\PassOptionsToPackage{hyphens}{url}
\PassOptionsToPackage{dvipsnames,svgnames,x11names}{xcolor}
%
\documentclass[
  letterpaper,
  DIV=11,
  numbers=noendperiod]{scrreprt}
\usepackage{xcolor}
\usepackage{amsmath,amssymb}
\setcounter{secnumdepth}{5}
\usepackage{iftex}
\ifPDFTeX
  \usepackage[T1]{fontenc}
  \usepackage[utf8]{inputenc}
  \usepackage{textcomp} % provide euro and other symbols
\else % if luatex or xetex
  \usepackage{unicode-math} % this also loads fontspec
  \defaultfontfeatures{Scale=MatchLowercase}
  \defaultfontfeatures[\rmfamily]{Ligatures=TeX,Scale=1}
\fi
\usepackage{lmodern}
\ifPDFTeX\else
  % xetex/luatex font selection
\fi
% Use upquote if available, for straight quotes in verbatim environments
\IfFileExists{upquote.sty}{\usepackage{upquote}}{}
\IfFileExists{microtype.sty}{% use microtype if available
  \usepackage[]{microtype}
  \UseMicrotypeSet[protrusion]{basicmath} % disable protrusion for tt fonts
}{}
\makeatletter
\@ifundefined{KOMAClassName}{% if non-KOMA class
  \IfFileExists{parskip.sty}{%
    \usepackage{parskip}
  }{% else
    \setlength{\parindent}{0pt}
    \setlength{\parskip}{6pt plus 2pt minus 1pt}}
}{% if KOMA class
  \KOMAoptions{parskip=half}}
\makeatother
% Make \paragraph and \subparagraph free-standing
\makeatletter
\ifx\paragraph\undefined\else
  \let\oldparagraph\paragraph
  \renewcommand{\paragraph}{
    \@ifstar
      \xxxParagraphStar
      \xxxParagraphNoStar
  }
  \newcommand{\xxxParagraphStar}[1]{\oldparagraph*{#1}\mbox{}}
  \newcommand{\xxxParagraphNoStar}[1]{\oldparagraph{#1}\mbox{}}
\fi
\ifx\subparagraph\undefined\else
  \let\oldsubparagraph\subparagraph
  \renewcommand{\subparagraph}{
    \@ifstar
      \xxxSubParagraphStar
      \xxxSubParagraphNoStar
  }
  \newcommand{\xxxSubParagraphStar}[1]{\oldsubparagraph*{#1}\mbox{}}
  \newcommand{\xxxSubParagraphNoStar}[1]{\oldsubparagraph{#1}\mbox{}}
\fi
\makeatother

\usepackage{color}
\usepackage{fancyvrb}
\newcommand{\VerbBar}{|}
\newcommand{\VERB}{\Verb[commandchars=\\\{\}]}
\DefineVerbatimEnvironment{Highlighting}{Verbatim}{commandchars=\\\{\}}
% Add ',fontsize=\small' for more characters per line
\usepackage{framed}
\definecolor{shadecolor}{RGB}{241,243,245}
\newenvironment{Shaded}{\begin{snugshade}}{\end{snugshade}}
\newcommand{\AlertTok}[1]{\textcolor[rgb]{0.68,0.00,0.00}{#1}}
\newcommand{\AnnotationTok}[1]{\textcolor[rgb]{0.37,0.37,0.37}{#1}}
\newcommand{\AttributeTok}[1]{\textcolor[rgb]{0.40,0.45,0.13}{#1}}
\newcommand{\BaseNTok}[1]{\textcolor[rgb]{0.68,0.00,0.00}{#1}}
\newcommand{\BuiltInTok}[1]{\textcolor[rgb]{0.00,0.23,0.31}{#1}}
\newcommand{\CharTok}[1]{\textcolor[rgb]{0.13,0.47,0.30}{#1}}
\newcommand{\CommentTok}[1]{\textcolor[rgb]{0.37,0.37,0.37}{#1}}
\newcommand{\CommentVarTok}[1]{\textcolor[rgb]{0.37,0.37,0.37}{\textit{#1}}}
\newcommand{\ConstantTok}[1]{\textcolor[rgb]{0.56,0.35,0.01}{#1}}
\newcommand{\ControlFlowTok}[1]{\textcolor[rgb]{0.00,0.23,0.31}{\textbf{#1}}}
\newcommand{\DataTypeTok}[1]{\textcolor[rgb]{0.68,0.00,0.00}{#1}}
\newcommand{\DecValTok}[1]{\textcolor[rgb]{0.68,0.00,0.00}{#1}}
\newcommand{\DocumentationTok}[1]{\textcolor[rgb]{0.37,0.37,0.37}{\textit{#1}}}
\newcommand{\ErrorTok}[1]{\textcolor[rgb]{0.68,0.00,0.00}{#1}}
\newcommand{\ExtensionTok}[1]{\textcolor[rgb]{0.00,0.23,0.31}{#1}}
\newcommand{\FloatTok}[1]{\textcolor[rgb]{0.68,0.00,0.00}{#1}}
\newcommand{\FunctionTok}[1]{\textcolor[rgb]{0.28,0.35,0.67}{#1}}
\newcommand{\ImportTok}[1]{\textcolor[rgb]{0.00,0.46,0.62}{#1}}
\newcommand{\InformationTok}[1]{\textcolor[rgb]{0.37,0.37,0.37}{#1}}
\newcommand{\KeywordTok}[1]{\textcolor[rgb]{0.00,0.23,0.31}{\textbf{#1}}}
\newcommand{\NormalTok}[1]{\textcolor[rgb]{0.00,0.23,0.31}{#1}}
\newcommand{\OperatorTok}[1]{\textcolor[rgb]{0.37,0.37,0.37}{#1}}
\newcommand{\OtherTok}[1]{\textcolor[rgb]{0.00,0.23,0.31}{#1}}
\newcommand{\PreprocessorTok}[1]{\textcolor[rgb]{0.68,0.00,0.00}{#1}}
\newcommand{\RegionMarkerTok}[1]{\textcolor[rgb]{0.00,0.23,0.31}{#1}}
\newcommand{\SpecialCharTok}[1]{\textcolor[rgb]{0.37,0.37,0.37}{#1}}
\newcommand{\SpecialStringTok}[1]{\textcolor[rgb]{0.13,0.47,0.30}{#1}}
\newcommand{\StringTok}[1]{\textcolor[rgb]{0.13,0.47,0.30}{#1}}
\newcommand{\VariableTok}[1]{\textcolor[rgb]{0.07,0.07,0.07}{#1}}
\newcommand{\VerbatimStringTok}[1]{\textcolor[rgb]{0.13,0.47,0.30}{#1}}
\newcommand{\WarningTok}[1]{\textcolor[rgb]{0.37,0.37,0.37}{\textit{#1}}}

\usepackage{longtable,booktabs,array}
\newcounter{none} % for unnumbered tables
\usepackage{calc} % for calculating minipage widths
% Correct order of tables after \paragraph or \subparagraph
\usepackage{etoolbox}
\makeatletter
\patchcmd\longtable{\par}{\if@noskipsec\mbox{}\fi\par}{}{}
\makeatother
% Allow footnotes in longtable head/foot
\IfFileExists{footnotehyper.sty}{\usepackage{footnotehyper}}{\usepackage{footnote}}
\makesavenoteenv{longtable}
\usepackage{graphicx}
\makeatletter
\newsavebox\pandoc@box
\newcommand*\pandocbounded[1]{% scales image to fit in text height/width
  \sbox\pandoc@box{#1}%
  \Gscale@div\@tempa{\textheight}{\dimexpr\ht\pandoc@box+\dp\pandoc@box\relax}%
  \Gscale@div\@tempb{\linewidth}{\wd\pandoc@box}%
  \ifdim\@tempb\p@<\@tempa\p@\let\@tempa\@tempb\fi% select the smaller of both
  \ifdim\@tempa\p@<\p@\scalebox{\@tempa}{\usebox\pandoc@box}%
  \else\usebox{\pandoc@box}%
  \fi%
}
% Set default figure placement to htbp
\def\fps@figure{htbp}
\makeatother





\setlength{\emergencystretch}{3em} % prevent overfull lines

\providecommand{\tightlist}{%
  \setlength{\itemsep}{0pt}\setlength{\parskip}{0pt}}



 


\KOMAoption{captions}{tableheading}
\makeatletter
\@ifpackageloaded{bookmark}{}{\usepackage{bookmark}}
\makeatother
\makeatletter
\@ifpackageloaded{caption}{}{\usepackage{caption}}
\AtBeginDocument{%
\ifdefined\contentsname
  \renewcommand*\contentsname{Table of contents}
\else
  \newcommand\contentsname{Table of contents}
\fi
\ifdefined\listfigurename
  \renewcommand*\listfigurename{List of Figures}
\else
  \newcommand\listfigurename{List of Figures}
\fi
\ifdefined\listtablename
  \renewcommand*\listtablename{List of Tables}
\else
  \newcommand\listtablename{List of Tables}
\fi
\ifdefined\figurename
  \renewcommand*\figurename{Figure}
\else
  \newcommand\figurename{Figure}
\fi
\ifdefined\tablename
  \renewcommand*\tablename{Table}
\else
  \newcommand\tablename{Table}
\fi
}
\@ifpackageloaded{float}{}{\usepackage{float}}
\floatstyle{ruled}
\@ifundefined{c@chapter}{\newfloat{codelisting}{h}{lop}}{\newfloat{codelisting}{h}{lop}[chapter]}
\floatname{codelisting}{Listing}
\newcommand*\listoflistings{\listof{codelisting}{List of Listings}}
\usepackage{amsthm}
\theoremstyle{remark}
\AtBeginDocument{\renewcommand*{\proofname}{Proof}}
\newtheorem*{remark}{Remark}
\newtheorem*{solution}{Solution}
\newtheorem{refremark}{Remark}[chapter]
\newtheorem{refsolution}{Solution}[chapter]
\makeatother
\makeatletter
\makeatother
\makeatletter
\@ifpackageloaded{caption}{}{\usepackage{caption}}
\@ifpackageloaded{subcaption}{}{\usepackage{subcaption}}
\makeatother
\usepackage{bookmark}
\IfFileExists{xurl.sty}{\usepackage{xurl}}{} % add URL line breaks if available
\urlstyle{same}
% fallback for those not using the hyperref driver hyperxmp:
\makeatletter
\@ifundefined{xmpquote}{\newcommand{\xmpquote}[1]{#1}}{}
\makeatother
\hypersetup{
  pdftitle={Should we trust academia?},
  pdfauthor={Aiden Ross},
  colorlinks=true,
  linkcolor={blue},
  filecolor={Maroon},
  citecolor={Blue},
  urlcolor={Blue},
  pdfcreator={LaTeX via pandoc}}


\title{Should we trust academia?}
\author{Aiden Ross}
\date{2026-11-05}
\begin{document}
\maketitle

\renewcommand*\contentsname{Table of contents}
{
\setcounter{tocdepth}{2}
\tableofcontents
}

\bookmarksetup{startatroot}

\chapter*{Preface}\label{preface}
\addcontentsline{toc}{chapter}{Preface}

\markboth{Preface}{Preface}

Augustus De Morgan wrote in \textbf{\emph{A Budget of Paradoxes, Volume
I}}, `Modern discoveries have not been made by large collections of
facts, with subsequent discussion, separation, and resulting deduction
of a truth thus rendered perceptible. A few facts have suggested an
hypothesis, which means a supposition, proper to explain them.' De
Morgan was a British mathematician famous for De Morgan's laws in logic
and set theory and for coining the term mathematical induction.

\begin{figure}

\centering{

\includegraphics[width=2.60417in,height=\textheight,keepaspectratio]{index_files/mediabag/500px-De_Morgan_Augu.jpg}

}

\caption{\label{fig-demorgan}Augustus De Morgan}

\end{figure}%

When he authored this book, he used an old meaning of the word
paradoxes. He explains this in the introduction: ``I use the word in the
old sense: a paradox is something which is apart from general opinion,
either in subject-matter, method, or conclusion.''

\bookmarksetup{startatroot}

\chapter{What is truth?}\label{what-is-truth}

Seems simple at first, but really, what is truth? An example of truth
you may hear from a sermon is that \(2+2=4\), but is that true? Well, it
is true that \(2+2=4\) in the decimal system, but nobody is forcing us
to use the decimal system. What if we use trinary or base-3? Base-3 is a
system where we group things by 3 instead of 10 like in our normal
decimal system. So if we have 2 things, that fits easily in our first
group, but if we add 2 more things, we have 1 group of 3 and 1 group of
1. Written in normal notation:

\[ 2 + 2 = 11 \]

It is usual to assume decimal notation, so often we prefer to write it
specifying the base as a subscript like this:

\[ 2_3 + 2_3 = 11_3 \]

That is an example of a system or model that we use to keep things
simple. Unfortunately, I am not able to define truth for you, as to this
day philosophers debate what truth is. Day-to-day, what tends to matter
is if something works reliably and this often means creating models that
are accurate for what we use them for but not truly representative of
the core of how something functions. As humans, we are limited in this
way: we have to define things to a limited precision.

\begin{figure}

\centering{

\includegraphics[width=4.16667in,height=\textheight,keepaspectratio]{index_files/mediabag/Two_hand,_ten_finger.jpg}

}

\caption{\label{fig-hands}10 fingers makes base-10 seem easy}

\end{figure}%

\bookmarksetup{startatroot}

\chapter{Dogmatism}\label{dogmatism}

\section{Useful definitions}\label{useful-definitions}

\begin{itemize}
\tightlist
\item
  Axiom: A claim that cannot be proven
\item
  Tautology: A statement that is always true
\item
  Contradiction: A statement that is always false
\end{itemize}

\section{Priors}\label{priors}

In my studies as an undergraduate mathematics major, one of the required
courses was \emph{Techniques of Proof}. This class aims to teach
students how to prove things absolutely. After learning formal logic, we
were introduced to a system of axioms. For logic to work, we have to
have some base assumptions. In our day-to-day lives, priors act as our
axioms or base assumptions.

Imagine you're in an elevator. In order to reach the 5th floor, you
likely know, based on past experiences, that you need to press the
button with the number 5 inscribed on it.

Consider the possible results from that button press:

{\def\LTcaptype{none} % do not increment counter
\begin{longtable}[]{@{}lcc@{}}
\toprule\noalign{}
& Button Functions Normally & Button Malfunctions \\
\midrule\noalign{}
\endhead
\bottomrule\noalign{}
\endlastfoot
result & ? & ? \\
\end{longtable}
}

Assuming you have mostly only interacted with functional elevators, then
you might rate your confidence in each outcome as:

{\def\LTcaptype{none} % do not increment counter
\begin{longtable}[]{@{}lcc@{}}
\toprule\noalign{}
& Button Functions Normally & Button Malfunctions \\
\midrule\noalign{}
\endhead
\bottomrule\noalign{}
\endlastfoot
result & \(90\%\) & \(10\%\) \\
\end{longtable}
}

These numbers are not measurements of how often the button works; they
are a rating of how sure you feel.

\section{Posteriors}\label{posteriors}

The posterior is your updated confidence after seeing some evidence.
Keeping the elevator example, imagine you are on the 3rd floor and the
light indicating someone is traveling upwards beeps on, and shortly
after, a person steps out onto the 3rd floor. That acts as evidence that
the elevator is currently functioning normally, and thus it would make
sense to change your estimates.

The weight or impact of the evidence here is not important, but you are
probably more certain now. This would create your posteriors, and they
might look like this:

{\def\LTcaptype{none} % do not increment counter
\begin{longtable}[]{@{}lcc@{}}
\toprule\noalign{}
& Button Functions Normally & Button Malfunctions \\
\midrule\noalign{}
\endhead
\bottomrule\noalign{}
\endlastfoot
result & \(95\%\) & \(5\%\) \\
\end{longtable}
}

Keep in mind these are not frequentist estimates. With these tables you
are not saying that in a series of elevator rides you always have a 95\%
chance of the button working normally; instead, these are subjective
ratings, like how you may tell your doctor that on a scale of one to ten
your pain is at a two.

\section{Dogma}\label{dogma}

Finally, let's create the dogmatic prior probability tables:

{\def\LTcaptype{none} % do not increment counter
\begin{longtable}[]{@{}lcc@{}}
\toprule\noalign{}
& Button Functions Normally & Button Malfunctions \\
\midrule\noalign{}
\endhead
\bottomrule\noalign{}
\endlastfoot
result & \(0\%\) & \(100\%\) \\
\end{longtable}
}

or

{\def\LTcaptype{none} % do not increment counter
\begin{longtable}[]{@{}lcc@{}}
\toprule\noalign{}
& Button Functions Normally & Button Malfunctions \\
\midrule\noalign{}
\endhead
\bottomrule\noalign{}
\endlastfoot
result & \(100\%\) & \(0\%\) \\
\end{longtable}
}

In both cases, one of the choices is set to 100\% certainty, which,
outside of a tautology or contradiction, is considered malformed by
most. The reason for this: what happens if the \(0\%\) case turns out to
be true? That would contradict the claim carried by \(0\%\) confidence,
which is that x can never happen. This would be fairly clear if you
shared the elevator with a person who fully claims that the button won't
work, even after the elevator successfully arrives.

Thus, the dogmatic posteriors are:

{\def\LTcaptype{none} % do not increment counter
\begin{longtable}[]{@{}lcc@{}}
\toprule\noalign{}
& Button Functions Normally & Button Malfunctions \\
\midrule\noalign{}
\endhead
\bottomrule\noalign{}
\endlastfoot
result & \(0\%\) & \(100\%\) \\
\end{longtable}
}

or

{\def\LTcaptype{none} % do not increment counter
\begin{longtable}[]{@{}lcc@{}}
\toprule\noalign{}
& Button Functions Normally & Button Malfunctions \\
\midrule\noalign{}
\endhead
\bottomrule\noalign{}
\endlastfoot
result & \(100\%\) & \(0\%\) \\
\end{longtable}
}

In dogmatism, the priors always equal the posteriors, regardless of
evidence.

\begin{Shaded}
\begin{Highlighting}[]
\NormalTok{viewof certain = Inputs.range([0, 100], \{label: html\textasciigrave{}certainty\textasciigrave{}, step: 0.1\})}

\NormalTok{dogma\_score = function(score)}
\NormalTok{\{}
\NormalTok{  if (score == 100 || score == 0)}
\NormalTok{    return "Dogmatic"}
\NormalTok{  else }
\NormalTok{    return "Not Dogmatic"}
\NormalTok{\}}

\NormalTok{data = [}
\NormalTok{\{outcome: "yes", percent: certain, dogma: dogma\_score(certain)\},}
\NormalTok{\{outcome: "no",  percent: 100 {-} certain, dogma: dogma\_score(100 {-} certain)\}}
\NormalTok{]}

\NormalTok{Plot.barY(data, \{x: "outcome", y: "percent", fill:"dogma"\}).plot(\{}
\NormalTok{  y: \{domain: [0, 100]\},}
\NormalTok{color: \{scheme: "BuRd", legend: true\}}
\NormalTok{  \})}
\end{Highlighting}
\end{Shaded}

\bookmarksetup{startatroot}

\chapter{Models and Systems}\label{models-and-systems}

Finally, before reaching the primary topic of this site we have one last
topic I would like to explain. In mathematics models have a formal
definition but it isn't really necessary to use it here. The important
thing is that a model uses assumptions and works on a given universe.

For example, consider the band Ghost. In 2013 they did a cover of a song
\emph{If You Have Ghost}. In the lyrics one of the lines reads almost
like an implication: ``If you have ghosts, you have everything''.
Normally you could read this meaning metaphorically but if we change the
model through definitions we can change the meaning.

Let ghosts mean specifically the band and let everything be defined as
music. This will be our definition for the model. Then consider that if
you listen to Ghost's cover that you have ghosts. Thus in our model by
modus ponens we have everything or in other words we have music.

For a more formal version of this quasi-proof:

\begin{proof}
~

Definitions for the model, followed by modus ponens:

\[
\begin{array}{l}
\text{ghosts} := \text{the band Ghost} \\
\text{everything} := \text{music} \\[1em]
\text{ghosts} \implies \text{everything} \\
\text{ghosts} \\
\hline
\therefore \text{everything}
\end{array}
\]
\end{proof}

Here our universe includes the band ghosts and music but if we used
other definitions that did not work that would make the statement
incorrect. For example, consider that as a human I am pretty sure that I
can not fly. So if I defined ghosts to mean human and music to mean
flight then the implication becomes false making it nonsensacle. The
important thing to note is legs and the ability to fly in a given moment
are not related for the most part so I could make the statement true by
just getting on a plain.

\bookmarksetup{startatroot}

\chapter{Pathos}\label{pathos}

\bookmarksetup{startatroot}

\chapter{Summary}\label{summary}

In summary, this book has no content whatsoever.

todo summary.

\bookmarksetup{startatroot}

\chapter*{References}\label{references}
\addcontentsline{toc}{chapter}{References}

\markboth{References}{References}

\protect\phantomsection\label{refs}




\end{document}
